\documentclass[11pt]{article}
\usepackage[margin=1in]{geometry}
\usepackage{amsmath}
\usepackage{amssymb}
\usepackage{tcolorbox}

\title{Problem 2b: Step-by-Step Calculation \\ (Getting $\lambda_1$)}
\author{ECO 310 - Detailed Walkthrough}
\date{February 23, 2026}

\begin{document}

\maketitle

\section{The Setup}

\textbf{The Lagrangian (Convention 2):}
\begin{equation}
\mathcal{L} = -p_1 x_1 - p_2 x_2 + \lambda_1 x_1 + \lambda_2 x_2 + \lambda_3(x_1 + x_2 - u^*)
\end{equation}

\textbf{The guess we're testing:}
\begin{itemize}
\item Utility constraint is binding: $x_1 + x_2 = u^*$
\item Nonnegativity on good 1 is binding: $x_1 = 0$
\item Nonnegativity on good 2 is non-binding: $x_2 > 0$
\end{itemize}

\textbf{What this means:}
\begin{itemize}
\item Since $x_1 = 0$, we're buying NO good 1
\item Since $x_2 > 0$, we're buying SOME good 2
\item Combined with $x_1 + x_2 = u^*$, we get $x_2 = u^*$ (buy only good 2)
\end{itemize}

This guess makes sense when \textbf{good 2 is cheaper}: $p_2 < p_1$

\section{The General Optimality Conditions}

From taking derivatives of the Lagrangian:

\textbf{FOC with respect to $x_1$:}
\begin{equation}
\frac{\partial \mathcal{L}}{\partial x_1} = -p_1 + \lambda_1 + \lambda_3 \leq 0, \quad x_1 \geq 0, \quad x_1(-p_1 + \lambda_1 + \lambda_3) = 0
\end{equation}

\textbf{FOC with respect to $x_2$:}
\begin{equation}
\frac{\partial \mathcal{L}}{\partial x_2} = -p_2 + \lambda_2 + \lambda_3 \leq 0, \quad x_2 \geq 0, \quad x_2(-p_2 + \lambda_2 + \lambda_3) = 0
\end{equation}

\textbf{FOC with respect to $\lambda_3$:}
\begin{equation}
\frac{\partial \mathcal{L}}{\partial \lambda_3} = x_1 + x_2 - u^* \geq 0, \quad \lambda_3 \geq 0, \quad \lambda_3(x_1 + x_2 - u^*) = 0
\end{equation}

\section{Step 1: Use Complementary Slackness for $x_2$}

Since $x_2 > 0$ (strictly positive), the complementary slackness condition says:
\begin{equation}
x_2 \cdot (-p_2 + \lambda_2 + \lambda_3) = 0
\end{equation}

Because $x_2 > 0$, we need:
\begin{equation}
-p_2 + \lambda_2 + \lambda_3 = 0
\end{equation}

The term in parentheses MUST equal zero (not just $\leq 0$).

\section{Step 2: Use Complementary Slackness for $\lambda_2$}

We also have complementary slackness for the constraint $x_2 \geq 0$:
\begin{equation}
\lambda_2 \cdot x_2 = 0
\end{equation}

Since $x_2 > 0$, we must have:
\begin{equation}
\lambda_2 = 0
\end{equation}

\section{Step 3: Solve for $\lambda_3$}

Substitute $\lambda_2 = 0$ into equation (6):
\begin{align}
-p_2 + 0 + \lambda_3 &= 0 \\
\lambda_3 &= p_2
\end{align}

\begin{tcolorbox}[colback=green!5!white,colframe=green!75!black,title=Key Result]
$$\lambda_3 = p_2$$

This is the marginal cost of utility. Since we're buying only good 2, each extra unit of utility costs $p_2$.
\end{tcolorbox}

\section{Step 4: Apply the $x_1$ Condition}

Since $x_1 = 0$ (binding constraint), we have the inequality:
\begin{equation}
-p_1 + \lambda_1 + \lambda_3 \leq 0
\end{equation}

Substitute $\lambda_3 = p_2$:
\begin{align}
-p_1 + \lambda_1 + p_2 &\leq 0 \\
\lambda_1 + p_2 &\leq p_1 \\
\lambda_1 &\leq p_1 - p_2
\end{align}

\section{Step 5: Determine $\lambda_1$}

At the optimal solution, this inequality typically \textbf{binds} (holds with equality):
\begin{equation}
\lambda_1 = p_1 - p_2
\end{equation}

\begin{tcolorbox}[colback=yellow!10!white,colframe=orange!75!black,title=Important!]
\textbf{The answer is:}
$$\lambda_1 = p_1 - p_2$$

\textbf{NOT} $\lambda_1 = p_2 - p_1$ (I had an error in the previous guide!)
\end{tcolorbox}

\section{Step 6: Check Validity}

For this solution to be valid, we need $\lambda_1 \geq 0$:
\begin{align}
\lambda_1 &\geq 0 \\
p_1 - p_2 &\geq 0 \\
p_1 &\geq p_2
\end{align}

\textbf{Wait, something seems wrong!}

If we're buying only good 2 ($x_1 = 0$, $x_2 = u^*$), this should happen when good 2 is \textit{cheaper}, meaning $p_2 < p_1$.

But our condition says $p_1 \geq p_2$, which is the same thing!

So actually, the solution IS valid when $p_1 > p_2$ (good 1 is more expensive than good 2).

\section{Concrete Example}

\textbf{Setup:}
\begin{itemize}
\item Good 1 (coffee) costs \$5: $p_1 = 5$
\item Good 2 (tea) costs \$3: $p_2 = 3$
\item Target utility: $u^* = 10$
\end{itemize}

\textbf{Guess:} Buy only good 2 (tea is cheaper)
\begin{itemize}
\item $x_1 = 0$ (no coffee)
\item $x_2 = 10$ (10 teas)
\end{itemize}

\textbf{Calculate shadow prices:}

From Step 3:
\begin{equation}
\lambda_3 = p_2 = 3
\end{equation}

From Step 5:
\begin{equation}
\lambda_1 = p_1 - p_2 = 5 - 3 = 2
\end{equation}

Also:
\begin{equation}
\lambda_2 = 0
\end{equation}

\textbf{Check validity:}
\begin{equation}
\lambda_1 = 2 > 0 \quad \checkmark
\end{equation}

\textbf{Economic interpretation:}
\begin{itemize}
\item $\lambda_1 = 2$: If you could buy "negative coffee" (get money back for not buying coffee), each unit would save you \$2

Actually, that doesn't quite make sense. Let me think about this differently...

\item $\lambda_1 = 2$ measures how much your cost would \textit{increase} if you were forced to buy 1 unit of good 1

\item Since coffee costs \$5 but tea costs \$3, replacing 1 tea with 1 coffee increases cost by $5 - 3 = 2$ dollars

\item $\lambda_3 = 3$: Marginal cost of utility is \$3 (the price of tea, which you're buying)
\end{itemize}

\section{Summary: Step-by-Step Process}

\begin{tcolorbox}[colback=blue!5!white,colframe=blue!75!black,title=How to Solve Problem 2b]

\textbf{Given the guess:} $x_1 = 0$, $x_2 > 0$, $x_1 + x_2 = u^*$

\textbf{Step 1:} Use complementary slackness
\begin{itemize}
\item Since $x_2 > 0$: the $x_2$ FOC holds with equality
\item Since $x_2 > 0$: we have $\lambda_2 = 0$
\end{itemize}

\textbf{Step 2:} Solve for $\lambda_3$
\begin{itemize}
\item From $x_2$ FOC: $-p_2 + \lambda_2 + \lambda_3 = 0$
\item Substitute $\lambda_2 = 0$: $\lambda_3 = p_2$
\end{itemize}

\textbf{Step 3:} Solve for $\lambda_1$
\begin{itemize}
\item From $x_1$ FOC: $-p_1 + \lambda_1 + \lambda_3 \leq 0$
\item Substitute $\lambda_3 = p_2$: $\lambda_1 \leq p_1 - p_2$
\item At optimum: $\lambda_1 = p_1 - p_2$
\end{itemize}

\textbf{Step 4:} Find $x_1$ and $x_2$
\begin{itemize}
\item From guess: $x_1 = 0$
\item From utility constraint: $x_2 = u^*$
\end{itemize}

\textbf{Step 5:} Check validity
\begin{itemize}
\item Need $\lambda_1 \geq 0$: requires $p_1 \geq p_2$
\item This makes sense: we buy only good 2 when it's cheaper
\end{itemize}

\end{tcolorbox}

\section{Why Does $-p_1 + \lambda_1 + \lambda_3 \leq 0$ Become $\lambda_1 = p_1 - p_2$?}

Let me break down this algebra very carefully:

\textbf{Starting equation:}
\begin{equation}
-p_1 + \lambda_1 + \lambda_3 \leq 0
\end{equation}

\textbf{Add $p_1$ to both sides:}
\begin{equation}
\lambda_1 + \lambda_3 \leq p_1
\end{equation}

\textbf{Subtract $\lambda_3$ from both sides:}
\begin{equation}
\lambda_1 \leq p_1 - \lambda_3
\end{equation}

\textbf{Substitute $\lambda_3 = p_2$:}
\begin{equation}
\lambda_1 \leq p_1 - p_2
\end{equation}

\textbf{At the optimum, assume this binds (equality):}
\begin{equation}
\lambda_1 = p_1 - p_2
\end{equation}

\begin{tcolorbox}[colback=green!5!white,colframe=green!75!black,title=Final Answer]
Given the guess that $x_1 = 0$, $x_2 = u^*$:

\begin{align}
x_1 &= 0 \\
x_2 &= u^* \\
\lambda_1 &= p_1 - p_2 \\
\lambda_2 &= 0 \\
\lambda_3 &= p_2
\end{align}

Valid when $p_1 \geq p_2$ (good 1 is at least as expensive as good 2)
\end{tcolorbox}

\end{document}
